\section{Conclusions}\label{sec:conclusions}

A 2.5D Fourier-reduced, quasistatic model of guided-wave propagation and
plate-edge reflection has been formulated for an n-type ZnO
piezoelectric-semiconductor plate with \textit{6mm} symmetry. The model retains
classical elasticity, piezoelectricity, dielectrics, strain-gradient elasticity,
electrostatics, carrier drift--diffusion, and the single additional
carrier--strain-gradient coupling
\[
\Delta H=\Delta n\,H_{jkl}\varepsilon_{jk,l}.
\]
No additional coupling, field, thermal effect, micro-inertia, quantum effect,
surface theory, or unrelated mechanism is introduced.

The numerical solution is a Legendre spectral thickness solver with an independent
edge-state formulation. The primary edge residual is \(3.00\times10^{-15}\) and the
control-volume power-balance error is \(7.31\times10^{-13}\). The primary isolated
backward self-power at \SI{2}{\micro\metre} is \SI{2.315}{\percent}; the
\(H/H^*=\pm0.05\) cases give \SI{2.286}{\percent} and \SI{2.266}{\percent}. These
are local isolated self-flux ratios, not total reflectance.

The manuscript does not claim a first-ever result relative to Qu and co-workers,
because the primary Appendix C equations remain inaccessible. It also does not claim
a full arbitrary 3D edge solution, broad-angle negative reflection, open-boundary
success, or calibrated physical \(H\) or flexoelectric coefficients. The main
scientific output is a mathematically consistent, reproducible, and explicitly
limited calculation of a solved backward channel in a specified piezoelectric-semiconductor
plate reduction.
